% §4 IoTBench: Benchmark Design — target 2 pages
% Purpose: sell the artifact.
% Key figure: Fig 1 — 5 architectural patterns side by side.
% Key table: Table 2 — scenarios × protocols × vulns × OWASP mapping.

\section{\systemname: A Network-Scale Benchmark}
\label{sec:iotbench}

% ── §4.1 design principles ────────────────────────────────────────────────
\subsection{Design principles}
\label{sec:iotbench:principles}

\systemname{} is built around four principles derived from the gap analysis in
\cref{sec:related}:

\begin{description}
  \item[Reproducibility.] The entire infrastructure deploys from Ansible
  playbooks against a Proxmox cluster; every vulnerability is injected
  idempotently with no manual post-deployment steps.
  \item[Architectural diversity.] Scenarios span five network patterns
  (\cref{fig:patterns}) with distinct attacker reachability graphs enforced by
  OpenWrt firewall rules, not just by service labels.
  \item[Protocol coverage.] Each scenario mixes multiple IoT and OT protocols
  (MQTT, Modbus, LoRaWAN gateways, CoAP, SNMP, plus standard SSH/HTTP/FTP/Telnet)
  to exercise protocol-specific reasoning.
  \item[Ground-truth granularity.] Every injected vulnerability is documented in
  YAML with \texttt{device}, \texttt{IP}, \texttt{role}, \texttt{severity},
  \texttt{category}, \texttt{OWASP IoT} and \texttt{MITRE ICS} mappings,
  optional \texttt{CVE}, \texttt{indicators} (expected detection strings), and
  a \texttt{verification} command.
\end{description}

% ── §4.2 five architectural patterns ──────────────────────────────────────
\subsection{Five architectural patterns}
\label{sec:iotbench:patterns}

\begin{figure*}[t]
  \centering
  % \includegraphics[width=\linewidth]{figures/fig1_patterns.pdf}
  \fbox{\parbox{0.95\linewidth}{\centering\textit{
    Fig.~1 placeholder: five topology diagrams side-by-side.\\
    (Flat)~— (Gateway-centric)~— (Segmented IT/OT)~— (Hub-star)~— (Edge-cloud)
  }}}
  \caption{The five architectural patterns in \systemname. Each pattern is enforced
  by OpenWrt firewall rules that constrain the attacker's reachability from the
  entry point, not by service labels alone.}
  \label{fig:patterns}
\end{figure*}

\todo{Write §4.2 prose — for each of the 5 patterns: (i)~one sentence naming the
pattern, (ii)~the real-world IoT deployment it models, (iii)~the firewall
reachability invariant.}

Concretely:
\begin{enumerate}
  \item \textbf{Flat.} All devices are directly reachable from the entry point;
  no pivot is required.
  \item \textbf{Gateway-centric.} A single IoT gateway mediates all access to
  back-end services; inter-service traffic is blocked without gateway compromise.
  \item \textbf{Segmented IT/OT.} Three zones (admin/IT/OT) with access-control
  lists; the OT zone (Modbus PLCs, HMI) is reachable only via a jump host.
  \item \textbf{Hub-centric star.} A central home-automation hub is the only
  device that can reach MQTT broker, DB, and camera; peripherals cannot talk
  to each other directly.
  \item \textbf{Edge-cloud pivot.} Edge subnet and cloud subnet are separated;
  only the edge gateway bridges them, making cloud-side services unreachable
  without edge compromise.
\end{enumerate}

% ── §4.3 vulnerability injection ──────────────────────────────────────────
\subsection{Deterministic vulnerability injection}
\label{sec:iotbench:inject}

Vulnerabilities are injected by a single Ansible role (\texttt{04\_inject\_vulns.yml})
parameterized by \texttt{scenario\_id}. Each scenario YAML lists \texttt{router\_vulns}
(OpenWrt-level: Telnet, WAN admin, FTP) and per-service \texttt{role} definitions
that trigger role-specific injections: \emph{e.g.}, the \texttt{mqtt\_broker} role
sets \texttt{allow\_anonymous true} and binds on \texttt{0.0.0.0:1883}; the
\texttt{ssh\_server} role installs weak credentials (\texttt{admin:admin}) and
enables \texttt{PasswordAuthentication}; the \texttt{db\_server} role starts
MariaDB with a passwordless \texttt{root}. Category distribution is shown in
\cref{tab:categories}.

\begin{table}[t]
\centering
\caption{Vulnerability categories and severity distribution across \systemname
(87~injected instances).}
\label{tab:categories}
\footnotesize
\begin{tabular}{@{}lrrrr@{}}
\toprule
Category & Count & CRIT & HIGH & MED/LOW \\
\midrule
\texttt{misconfiguration}      & 17 & 3 & 8 & 6 \\
\texttt{data\_exposure}        & 12 & 0 & 7 & 5 \\
\texttt{default\_credentials}  & 11 & 6 & 5 & 0 \\
\texttt{info\_disclosure}      & 10 & 0 & 0 & 10 \\
\texttt{no\_authentication}    &  9 & 2 & 4 & 3 \\
\texttt{insecure\_update}      &  4 & 0 & 0 & 4 \\
\texttt{weak\_crypto}          &  3 & 0 & 0 & 3 \\
\texttt{missing\_header}       &  3 & 0 & 0 & 3 \\
\texttt{known\_cve}            &  3 & 0 & 3 & 0 \\
\texttt{code\_injection}       &  3 & 2 & 1 & 0 \\
\texttt{privilege\_escalation} &  2 & 0 & 2 & 0 \\
\midrule
\textbf{Total}                 &\textbf{87} & 13 & 30 & 44 \\
\bottomrule
\end{tabular}
\end{table}

% ── §4.4 ground truth schema ──────────────────────────────────────────────
\subsection{Ground-truth schema}
\label{sec:iotbench:gt}

Each scenario ships a YAML file describing every injected vulnerability. An
entry example (from \texttt{scenario\_1.yaml}):

\begin{figure}[t]
\footnotesize
\begin{verbatim}
- id: V1
  device: s1-mqtt
  ip: "192.168.100.11"
  role: mqtt_broker
  title: "MQTT without authentication"
  severity: high
  category: misconfiguration
  owasp_iot: "I1 / I9"
  mitre_ics: "Initial Access, Collection"
  cve: null
  indicators:
    - "allow_anonymous true"
    - "Anonymous clients connected"
  verification: "mosquitto_sub -h 192.168.100.11 -t '#' -v"
\end{verbatim}
\caption{Ground-truth entry example.}
\label{lst:gt}
\end{figure}

% ── §4.5 evaluator ────────────────────────────────────────────────────────
\subsection{Evaluator}
\label{sec:iotbench:eval}

The evaluator (\S\,\ref{sec:harness:report} integrates results into Phase~5)
matches each LLM finding against ground truth in four passes, in order:
(1)~exact CVE match; (2)~\texttt{(ip, vuln\_type)}; (3)~\texttt{(ip, category)};
(4)~\texttt{(ip, severity)}. Severity weights are
$w = \{\text{critical}{=}4,\text{high}{=}3,\text{medium}{=}2,\text{low}{=}1\}$;
loose category matches apply $0.5\times$, severity mismatches $0.75\times$.
The composite score is the weighted True-Positive sum normalized by
$\sum_{g \in \text{GT}} w(g)$.

\paragraph{Evidence levels}
Beyond match accuracy, the evaluator tracks, per True Positive, the deepest
\emph{evidence level} achieved by the harness: \textsf{L1}~(detection — banner
or port), \textsf{L2}~(exploitation — authenticated action), \textsf{L3}~(data
exfiltration — sensitive content retrieved). We report
$\text{L\textit{k}-coverage} = |\{t \in \text{TP} : \text{level}(t) \geq k\}| / |\text{TP}|$.
This metric cannot be computed on single-host binary-flag benchmarks and
captures a dimension of pentest quality that existing LLM-pentest evaluations
conflate.
